\section{Программная реализация задачи}

\subsection{Структура программы}

Программа реализована на языке Java 17 с графическим интерфейсом на Swing. Структура проекта:

\begin{itemize}
\item \texttt{engine/} — численные методы интерполяции (чистые вычисления, без интерфейса)
\item \texttt{model/} — модели данных (InterpolationResult, InputData)
\item \texttt{service/} — загрузка данных (из таблицы, файла, функции)
\item \texttt{gui/} — графический интерфейс (Swing)
\end{itemize}

\subsection{Конечные разности (FiniteDifferenceTable.java)}

\begin{lstlisting}[style=javaStyle, label=code:fd]
public class FiniteDifferenceTable {

    public static double[][] build(double[] y) {
        int n = y.length;
        double[][] table = new double[n][n];
        for (int i = 0; i < n; i++) {
            table[i][0] = y[i];
        }
        for (int j = 1; j < n; j++) {
            for (int i = 0; i < n - j; i++) {
                table[i][j] = table[i + 1][j - 1] - table[i][j - 1];
            }
        }
        return table;
    }
}
\end{lstlisting}

\subsection{Полином Лагранжа (LagrangeSolver.java)}

\begin{lstlisting}[style=javaStyle, label=code:lagrange]
public class LagrangeSolver {

    public static InterpolationResult interpolate(double[] x, double[] y,
                                                  double xTarget) {
        int n = x.length;
        double result = 0.0;

        for (int i = 0; i < n; i++) {
            double term = y[i];
            for (int j = 0; j < n; j++) {
                if (i != j) {
                    term *= (xTarget - x[j]) / (x[i] - x[j]);
                }
            }
            result += term;
        }

        double[][] diffTable = FiniteDifferenceTable.build(y);
        return new InterpolationResult(
                "Lagrange Polynomial", xTarget, result, diffTable, n - 1);
    }
}
\end{lstlisting}

\subsection{Полином Ньютона с разделёнными разностями (NewtonSolver.java)}

\begin{lstlisting}[style=javaStyle, label=code:newton]
public class NewtonSolver {

    private static double[][] buildDividedDiffTable(double[] x, double[] y) {
        int n = y.length;
        double[][] table = new double[n][n];
        for (int i = 0; i < n; i++) {
            table[i][0] = y[i];
        }
        for (int j = 1; j < n; j++) {
            for (int i = 0; i < n - j; i++) {
                table[i][j] = (table[i + 1][j - 1] - table[i][j - 1])
                        / (x[i + j] - x[i]);
            }
        }
        return table;
    }

    /** Newton 1st interpolation formula (forward) */
    public static InterpolationResult interpolateForward(double[] x, double[] y,
                                                         double xTarget) {
        double[][] dd = buildDividedDiffTable(x, y);
        int n = x.length;
        double result = dd[0][0];
        double product = 1.0;

        for (int i = 1; i < n; i++) {
            product *= (xTarget - x[i - 1]);
            result += dd[0][i] * product;
        }
        // ...
        return new InterpolationResult(
                "Newton Divided Diff (Forward)", xTarget, result, fd, n - 1);
    }

    /** Newton 2nd interpolation formula (backward) */
    public static InterpolationResult interpolateBackward(double[] x, double[] y,
                                                          double xTarget) {
        double[][] dd = buildDividedDiffTable(x, y);
        int n = x.length;
        double result = dd[n - 1][0];
        double product = 1.0;

        for (int i = 1; i < n; i++) {
            product *= (xTarget - x[n - i]);
            result += dd[n - 1 - i][i] * product;
        }
        // ...
        return new InterpolationResult(
                "Newton Divided Diff (Backward)", xTarget, result, fd, n - 1);
    }
}
\end{lstlisting}

\subsection{Полином Гаусса (GaussSolver.java)}

\begin{lstlisting}[style=javaStyle, label=code:gauss]
public class GaussSolver {

    /** Gauss 1st formula (forward) */
    public static InterpolationResult interpolateForward(double[] x, double[] y,
                                                         double xTarget) {
        int n = y.length;
        double[][] ft = FiniteDifferenceTable.build(y);
        double h = x[1] - x[0];
        int c = n / 2;  // central index
        double q = (xTarget - x[c]) / h;
        double result = ft[c][0];

        for (int k = 1; k < n; k++) {
            double coeff = gaussForwardCoeff(q, k);
            int row = (k % 2 == 1) ? c - (k - 1) / 2 : c - k / 2;
            if (row >= 0 && row < n - k) {
                result += coeff * ft[row][k];
            }
        }
        // ...
    }

    /** Gauss 2nd formula (backward) */
    public static InterpolationResult interpolateBackward(double[] x, double[] y,
                                                          double xTarget) {
        // ... uses gaussBackwardCoeff with row = c - (k+1)/2 or c - k/2
    }
}
\end{lstlisting}

\subsection{Формула Стирлинга (StirlingSolver.java)}

Формула Стирлинга реализована как среднее арифметическое первой и второй формул Гаусса:

\begin{lstlisting}[style=javaStyle, label=code:stirling]
public class StirlingSolver {

    public static InterpolationResult interpolate(double[] x, double[] y,
                                                  double xTarget) {
        // ...
        double gf = gaussForwardEval(ft, c, q, n);
        double gb = gaussBackwardEval(ft, c, q, n);
        double result = (gf + gb) / 2.0;
        // ...
    }
}
\end{lstlisting}

\subsection{Формула Бесселя (BesselSolver.java)}

\begin{lstlisting}[style=javaStyle, label=code:bessel]
public class BesselSolver {

    public static InterpolationResult interpolate(double[] x, double[] y,
                                                  double xTarget) {
        // ...
        for (int k = 1; k < n; k++) {
            double coeff = besselCoeff(q, k);
            double diffVal;
            if (k % 2 == 0) {
                // even k: average of two adjacent even differences
                int row1 = c - k / 2;
                int row2 = c - k / 2 + 1;
                diffVal = (ft[row1][k] + ft[row2][k]) / 2.0;
            } else {
                // odd k: single odd difference
                int row = c - (k - 1) / 2;
                diffVal = ft[row][k];
            }
            result += coeff * diffVal;
        }
        // ...
    }
}
\end{lstlisting}

\subsection{Графический интерфейс (MainFrame.java)}

Приложение реализует три способа ввода данных:
\begin{enumerate}
\item \textbf{Ручной ввод} — ввод значений $x$ и $y$ через запятую
\item \textbf{Загрузка из файла} — чтение табличных данных из текстового файла
\item \textbf{Генерация по функции} — выбор функции ($\sin x$, $\cos x$, $e^x$, $x^2$, $\ln x$), интервала и количества точек
\end{enumerate}

Пользователь может выбрать произвольный набор методов интерполяции (Лагранж, Ньютон, Гаусс, Стирлинг, Бессель) и указать два значения аргумента $X_1$ и $X_2$.

Результаты отображаются в виде:
\begin{itemize}
\item Таблицы конечных разностей
\item Вычисленных значений для каждого метода
\item Графика с узлами интерполяции, интерполяционной кривой и отмеченными точками $X_1$, $X_2$
\end{itemize}

\section{Пример работы программы}

\begin{figure}[H]
    \centering
    \includegraphics[width=1\linewidth]{pic/screenshot_main.png}
    \caption{Основное окно программы с данными варианта 14}
    \label{fig:main}
\end{figure}

\begin{figure}[H]
    \centering
    \includegraphics[width=1\linewidth]{pic/screenshot_function.png}
    \caption{Интерполяция функции $\sin(x)$}
    \label{fig:func}
\end{figure}

\section{Анализ результатов}

\subsection{Проверка корректности}

Все методы интерполяции (Лагранж, Ньютон, Гаусс, Стирлинг, Бессель) дают одинаковые значения функции для одних и тех же узлов и аргументов. Это объясняется тем, что интерполяционный полином степени $n-1$ для $n$ узлов единственен — различные формулы представляют собой лишь разные формы записи одного и того же полинома.

\subsection{Сравнение результатов}

Для варианта 14 (таблица 1.4) получены следующие результаты:
\begin{itemize}
\item $f(1{,}112) \approx 0{,}749956$ (все методы)
\item $f(1{,}319) \approx 2{,}845036$ (все методы)
\end{itemize}

\subsection{Тестирование на некорректных данных}

Программа корректно обрабатывает:
\begin{itemize}
\item Несовпадение количества значений $x$ и $y$
\item Немонотонные значения $x$
\item Пустые данные
\item Некорректные числовые форматы
\end{itemize}

\subsection{Выводы}

\begin{enumerate}
\item Все реализованные методы интерполяции дают идентичные результаты при использовании полного набора конечных разностей
\item Для интерполяции вблизи начала таблицы целесообразно использовать первую формулу Ньютона
\item Для интерполяции вблизи середины таблицы удобны формулы Гаусса, Стирлинга и Бесселя
\item Формула Стирлинга даёт наилучшие результаты при $|q| \leq 0{,}25$
\item Формула Бесселя оптимальна при $|q| \approx 0{,}5$
\item Программная реализация с графическим интерфейсом позволяет наглядно сравнить результаты различных методов
\end{enumerate}

\section{Вывод}

В ходе выполнения лабораторной работы были изучены и реализованы основные методы интерполяции функций: полином Лагранжа, полином Ньютона с разделёнными разностями, полином Гаусса (первая и вторая формулы), а также дополнительные методы — формулы Стирлинга и Бесселя.

Вычислительная часть работы выполнена для варианта 14 (таблица 1.4): построена таблица конечных разностей, вычислены значения функции для $X_1 = 1{,}112$ и $X_2 = 1{,}319$ с использованием формул Ньютона, Гаусса и Лагранжа.

Программная часть реализована на языке Java с графическим интерфейсом Swing. Приложение поддерживает три способа ввода данных, позволяет выбрать набор методов интерполяции, отображает таблицу конечных разностей и строит график интерполяции. Программа протестирована на различных наборах данных, включая некорректные.

\end{enumerate}
